\documentclass[../../../include/open-logic-section]{subfiles}

\begin{document}

\olfileid{sth}{card-arithmetic}{opps}
\olsection{د بنسټيزو عملياتو تعريف}

څرنګه چې دلته د ترتيب ساتلو ته اړتيا نه لرو، د اصلي عددونو حساب
د ترتيبي عددونو تر حسابه يو څه په اسانه تعريفېږي۔ جمع، ضرب او
توان به يو ځاے تعريف کړو۔

\begin{defn}
کله چې $\cardfont{a}$ او $\cardfont{b}$ اصلي عددونه وي:
\begin{align*}
	\cardfont{a} \cardplus \cardfont{b} &\defis
	\card{\cardfont{a} \disjointsum \cardfont{b}}\\
	\cardfont{a} \cardtimes \cardfont{b} &\defis
	\card{\cardfont{a} \times \cardfont{b}}\\
	\cardexpo{\cardfont{a}}{\cardfont{b}} &\defis
	\card{\funfromto{\cardfont{b}}{\cardfont{a}}}
\end{align*}
چې دلته $\funfromto{X}{Y} = \Setabs{f}{f \text{تابع ده } X \to Y}$۔
(دا ښودل اسانه دي چې $\funfromto{X}{Y}$ د هر دوو سټونو $X$
او $Y$ دپاره شته؛ ثبوت يې د تمرين په توګه پرېږدو۔)
\end{defn}

\begin{prob}
په $\Zminus$ کښې ثابته کړئ چې $\funfromto{X}{Y}$ د هر دوو
سټونو $X$ او~$Y$ دپاره شته۔ په $\ZF$ کښې کار وکړئ او
$\setrank{\funfromto{X}{Y}}$ د $\setrank{X}$ او $\setrank{Y}$
له مخې محاسبه کړئ، لکه څنګه چې په
\olref[sth][ord-arithmetic][using-addition]{rankcomputation} کښې شوي دي۔
\end{prob}

د دې تعريف تشريح به مرسته وکړي۔ د جمعې په برخه کښې: دلته د
بېلې مجموعې مفهوم، $\disjointsum$، کارول کېږي، چې په
\olref[ord-arithmetic][add]{defdissum} کښې تعريف شوے دے؛ او
په اسانه ښکاري چې دا تعريف د متناهي حالتونو دپاره سمه پايله
ورکوي۔ د ضرب په برخه کښې: \olref[sfr][set][pai]{cardnmprod}
موږ ته وايي چې که $A$ د $n$ غړي او $B$ د $m$ غړي ولري، نو
$A \times B$ د $n \cdot m$ غړي لري؛ نو زموږ تعريف همدا تصور
له متناهي څخه نامتناهي ضرب ته غځوي۔ توان همداسې دے: د متناهي
حالت تصور نامتناهي حالت ته غځوو۔ په رښتيا، له ځينو اړخونو د
نامتناهي اصلي عددونو حساب د نامتناهي ترتيبي عددونو تر حسابه
«عادي» حساب ته ډېر ورته دے:

\begin{prop}\ollabel{cardplustimescommute}
$\cardplus$ او $\cardtimes$ دواړه تبادلي او اتحادي دي۔
\end{prop}

\begin{proof}
د تبادلي والي دپاره، د
\olref[cardinals][cardsasords]{lem:CardinalsBehaveRight} له مخې
بس ده پام وکړو چې $\cardeq{(\cardfont{a} \disjointsum
\cardfont{b})}{(\cardfont{b} \disjointsum \cardfont{a})}$ او
$\cardeq{(\cardfont{a} \times \cardfont{b})}{(\cardfont{b} \times
\cardfont{a})}$۔ اتحادي والى د تمرين په توګه پرېږدو۔
\end{proof}

\begin{prob}
ثابته کړئ چې $\cardplus$ او $\cardtimes$ اتحادي دي۔
\end{prob}

\begin{prop}
$A$ هغه وخت او يوازې هغه وخت نامتناهي دے چې
$\card{A} \cardplus 1 = 1 \cardplus \card{A} = \card{A}$ وي۔
\end{prop}

\begin{proof}
لکه په
\olref[cardinals][classing]{generalinfinitycharacter} کښې،
دا له \olref[ord-arithmetic][using-addition]{ordinfinitycharacter} او
\olref[cardinals][cardsasords]{lem:CardinalsBehaveRight} څخه
راځي۔
\end{proof}

دا څرګندوي چې ولې د ترتيبي او اصلي عددونو د جمعې او ضرب دپاره
بېلابېلو نښو ته اړتيا لرو: دا په رښتيا \emph{بېلابېل} عمليات
دي۔ د پايلو راتلونکې جوړه ښيي چې ترتيبي او اصلي توان هم
بېلابېل عمليات دي۔ (په ياد ولرئ چې له
\olref[z][infinity-again]{defnomega} څخه $2 = \{0,
1\}$ راځي):

\begin{lem}\ollabel{lem:SizePowerset2Exp}
$\card{\Pow{A}} = \cardexpo{2}{\card{A}}$، د هر $A$ دپاره۔
\end{lem}

\begin{proof}
د هر فرعي سټ $B \subseteq A$ دپاره، $\chi_B \in
\funfromto{A}{2}$ داسې وټاکئ:
\begin{align*}
	\chi_{B}(x) &\defis
	\begin{cases}
		1 & \text{که }x\in B\\
		0 & \text{کنه.}
	\end{cases}
\end{align*}
اوس $f(B) = \chi_B$ وټاکئ؛ دا يوه !!a{bijection}
$f \colon \Pow{A} \to \funfromto{A}{2}$ تعريفوي۔ نو
$\cardeq{\Pow{A}}{\funfromto{A}{2}}$۔ له دې
$\cardeq{\Pow{A}}{\funfromto{\card{A}}{2}}$ هم راځي، او بيا
$\card{\Pow{A}} = \card{\funfromto{\card{A}}{2}} =
2^{\card{A}}$۔
\end{proof}

دا لنډ ثبوت په اصل کښې د \olref[sfr][siz][red-alt]{sec} بحث
هم رانغاړي۔ هلته مو وښودله چې څنګه د $\Pow{\omega}$
نۀ‌شمېر وړتيا د نامتناهي دوه‌ګونو لړيو د سټ، $\Bin^\omega$،
نۀ‌شمېر وړتيا ته «راوګرځوو»۔ په عمل کښې، $\Bin^{\omega}$
هماغه $\funfromto{\omega}{2}$ دے؛ او مخکيني ثبوت وښودله چې
په \olref[sfr][siz][red-alt]{sec} کښې زموږ استدلال د
هر سټ~$A$ دپاره، د~$\omega$ پر ځاے، هم کار کوي۔ دا پايله د
اصلي توان په اړه يو ژر تر لاسه کېدونکے حقيقت هم راکوي:

\begin{cor}\ollabel{cantorcor}
$\cardfont{a} < \cardexpo{2}{\cardfont{a}}$ د هر اصلي عدد~$\cardfont{a}$ دپاره۔
\end{cor}

\begin{proof}
دا د کانتور له قضيې (\olref[sfr][siz][car]{thm:cantor}) او
\olref{lem:SizePowerset2Exp} څخه راځي۔
\end{proof}
\noindent
نو $\omega < \cardexpo{2}{\omega}$۔ خو پام وکړئ: دا پايله
د \emph{اصلي عددونو} د توان په اړه ده۔ له \emph{ترتيبي
عددونو} د توان سره بايد پرتله شي، ځکه په وروستي حالت کښې
$\omega = \ordexpo{2}{\omega}$ (وګورئ
\olref[ord-arithmetic][expo]{sec})۔

څرنګه چې اوس د اصلي توان په اړه غږېږو، د هغې «لارې» په اړه هم
لږ ډېر کره کېدے شو چې $\Real$ پرې !!{nonenumerable} دے۔

\begin{thm}\ollabel{continuumis2aleph0}
$\card{\Real} = \cardexpo{2}{\omega}$
\end{thm}

\begin{proof}[د ثبوت چوکاټ]
د دې د ثبوت ډېرې لارې شته۔ تر ټولو نېغه لاره دا ده چې
$\cardle{\Pow{\omega}}{\Real}$ او
$\cardle{\Real}{\Pow{\omega}}$ وښيو، بيا د
شرودر--برنشتاين له قضيې څخه $\cardeq{\Real}{\Pow{\omega}}$
واخلو، او له \olref[card-arithmetic][opps]{lem:SizePowerset2Exp}
څخه $\card{\Real} = \cardexpo{2}{\omega}$ ترلاسه کړو۔ د
$f \colon \Pow{\omega} \to \Real$ او
$g \colon \Real \to \Pow{\omega}$ غوندې يو پر يو تابعو
تعريف د (روښانوونکي) تمرين په توګه پرېږدو۔
\end{proof}

\begin{prob}
د \olref[sth][card-arithmetic][opps]{continuumis2aleph0}
ثبوت بشپړ کړئ؛ يعنې وښيئ چې $\cardle{\Pow{\omega}}{\Real}$
او $\cardle{\Real}{\Pow{\omega}}$۔
\end{prob}

\end{document}
